\documentclass[11pt]{article}

\usepackage[margin=0.8in]{geometry}
\usepackage{amsmath,amssymb}
\usepackage{enumitem}
\usepackage{xcolor}
\usepackage{tcolorbox}
\usepackage{fancyhdr}
\usepackage{lastpage}
\usepackage{hyperref}
\usepackage{microtype}
\usepackage{array}

\definecolor{uwpurple}{HTML}{4B2E83}
\definecolor{uwgold}{HTML}{B7A57A}
\definecolor{softpurple}{HTML}{F4F1FA}
\definecolor{softgold}{HTML}{FBF7EA}
\definecolor{deepgold}{HTML}{8A6F2A}

\tcbuselibrary{breakable,skins}

\pagestyle{fancy}
\fancyhf{}
\lhead{\textbf{MATH 394: Probability I}}
\rhead{\textbf{Makeup Midterm}}
\cfoot{\thepage\ of \pageref{LastPage}}
\setlength{\headheight}{14pt}

\setlength{\parindent}{0pt}
\setlength{\parskip}{0.45em}

\newcommand{\course}{MATH 394: Probability I}
\newcommand{\instructor}{Arman Jahangiri}
\newcommand{\department}{Department of Mathematics}
\newcommand{\university}{University of Washington}

\newtcolorbox{problemheader}{
    colback=softpurple,
    colframe=uwpurple,
    boxrule=0.8pt,
    arc=2mm,
    left=2mm,
    right=2mm,
    top=1mm,
    bottom=1mm
}



\newtcolorbox{policybox}{
    breakable,
    colback=softgold,
    colframe=deepgold,
    boxrule=0.8pt,
    arc=2mm,
    left=2.5mm,
    right=2.5mm,
    top=1.5mm,
    bottom=1.5mm
}

\begin{document}
\begin{titlepage}
\centering
\vspace*{0.05cm}

{\Large \textbf{\university}}\\
{\large \department}

\vspace{0.28cm}

{\LARGE \color{uwpurple}\textbf{Makeup Midterm Examination}}\\[0.15cm]
{\Large \textbf{\course}}\\[0.15cm]
{\large Summer 2026 \quad | \quad Instructor: \instructor}

\vspace{0.22cm}

% ------------------------------------------------------------
% EXAM INFORMATION
% ------------------------------------------------------------

\begin{tcolorbox}[
    width=0.94\textwidth,
    colback=softpurple,
    colframe=uwpurple,
    boxrule=0.9pt,
    arc=3mm
]
\centering
\renewcommand{\arraystretch}{1.15}
\begin{tabular}{>{\bfseries}p{0.19\textwidth}p{0.68\textwidth}}
Date: & \rule{0.35\textwidth}{0.4pt} \\
Answering time: & 60 minutes \\
Total points: & 55 \\
\end{tabular}
\end{tcolorbox}

\vspace{0.18cm}

% ------------------------------------------------------------
% STUDENT INFORMATION
% ------------------------------------------------------------

\begin{tcolorbox}[
    width=0.94\textwidth,
    colback=white,
    colframe=uwgold,
    boxrule=0.8pt,
    arc=3mm
]
\renewcommand{\arraystretch}{1.35}
\begin{tabular}{>{\bfseries}p{0.22\textwidth}p{0.67\textwidth}}
Name: & \hrulefill \\
Student ID: & \hrulefill \\
Signature: & \hrulefill \\
\end{tabular}
\end{tcolorbox}

\vspace{0.15cm}

% ------------------------------------------------------------
% EXAM POLICIES
% ------------------------------------------------------------

\begin{policybox}
\small

\textbf{Instructions}

\begin{itemize}[leftmargin=*,itemsep=0.05em,topsep=0.15em,parsep=0pt]
    \item This examination contains \textbf{five problems}. Verify that your copy is complete.
    \item Show all work and justify your answers. Unsupported answers may receive little or no credit.
    \item Clearly define any events or random variables that you introduce.
    \item Give exact answers unless an approximation is requested.
    \item Do not remove pages from the examination booklet.
\end{itemize}

\medskip

\textbf{Permitted materials}

The examination is closed-book and closed-notes. Students may use:

\begin{itemize}[leftmargin=*,itemsep=0.05em,topsep=0.12em,parsep=0pt]
    \item one handwritten 8.5 $\times$ 11 inch reference sheet (front and back),
    \item a non-programmable, non-graphing calculator capable of standard
    exponential and logarithmic computations (e.g.\ TI-30X IIS), and
    \item a standard normal table.
\end{itemize}

No other electronic devices, calculators, communication devices, or
external resources may be used during the examination.

\end{policybox}

\vspace{0.10cm}

% ------------------------------------------------------------
% SCORE TABLE
% ------------------------------------------------------------

\renewcommand{\arraystretch}{1.18}
\begin{tabular}{
|>{\centering\arraybackslash}p{0.16\textwidth}|
 >{\centering\arraybackslash}p{0.16\textwidth}|
}
\hline
\textbf{Problem} & \textbf{Score} \\
\hline
1 & /10 \\
\hline
2 & /10 \\
\hline
3 & /10 \\
\hline
4 & /15 \\
\hline
5 & /15 \\
\hline
\textbf{Total} & \textbf{/60} \\
\hline
\end{tabular}

\vspace{0.08cm}

{\footnotesize
By signing above, you affirm that the work is your own and that you followed
the examination rules.
}

\end{titlepage}
% ============================================================
% PROBLEM 1
% ============================================================

\begin{problemheader}
\textbf{Problem 1.} \hfill \textbf{10 points}
\end{problemheader}

A box contains 6 red balls, 5 blue balls, and 4 green balls. Five balls
are selected uniformly at random without replacement.

Find the probability that the selection contains exactly 2 red balls,
exactly 2 blue balls, and exactly 1 green ball.

Give an exact answer. A correct expression involving binomial coefficients
is acceptable.

\vspace{7cm}


\newpage

% ============================================================
% PROBLEM 2
% ============================================================

\begin{problemheader}
\textbf{Problem 2.} \hfill \textbf{10 points}
\end{problemheader}

A company purchases computer chips from three suppliers.

Supplier $A$ provides $50\%$ of the chips, Supplier $B$ provides $30\%$,
and Supplier $C$ provides the remaining $20\%$.

The probabilities that a chip is defective are
\[
P(D\mid A)=0.01,\qquad
P(D\mid B)=0.03,\qquad
P(D\mid C)=0.06,
\]
where $D$ denotes the event that a chip is defective.

A chip is selected at random and is found to be defective.

\begin{enumerate}[label=(\alph*),leftmargin=*,itemsep=0.8em]

    \item Find the probability that the defective chip came from
    Supplier $B$.
    \hfill \textbf{(6 points)}

    \item Find the probability that the defective chip came from
    either Supplier $B$ or Supplier $C$.
    \hfill \textbf{(4 points)}

\end{enumerate}

\vspace{7cm}


\newpage

% ============================================================
% PROBLEM 3
% ============================================================

\begin{problemheader}
\textbf{Problem 3.} \hfill \textbf{10 points}
\end{problemheader}

Let $A$ and $B$ be events.

\begin{enumerate}[label=(\alph*),leftmargin=*,itemsep=1em]

    \item Suppose
    \[
    P(A)=0.7,
    \qquad
    P(B)=0.6.
    \]
    What is the smallest possible value of $P(A\cap B)$?
    Justify your answer.
    \hfill \textbf{(3 points)}

    \item Suppose $A$ and $B$ are disjoint and
    \[
    P(A)>0,
    \qquad
    P(B)>0.
    \]
    Can $A$ and $B$ be independent? Justify your answer using the
    definition of independence.
    \hfill \textbf{(3 points)}

    \item Suppose $A$ and $B$ are independent. Prove that
    $A$ and $B^c$ are also independent.
    \hfill \textbf{(4 points)}

\end{enumerate}

\vspace{6cm}


\newpage

% ============================================================
% PROBLEM 4
% ============================================================

\begin{problemheader}
\textbf{Problem 4.} \hfill \textbf{15 points}
\end{problemheader}

Suppose that $X$ has probability density function
\[
f_X(x)=
\begin{cases}
c(1+x), & 0<x<2,\\
0, & \text{otherwise},
\end{cases}
\]
where $c$ is a constant.

\begin{enumerate}[label=(\alph*),leftmargin=*,itemsep=0.8em]

    \item Find $c$.
    \hfill \textbf{(3 points)}

    \item Find the cumulative distribution function $F_X(x)$
    for all real $x$.
    \hfill \textbf{(5 points)}

    \item Find
    \[
    P\left(\frac12<X<\frac32\right).
    \]
    \hfill \textbf{(2 points)}

    \item Find $\mathbb{E}[X]$ and $\operatorname{Var}(X)$.
    \hfill \textbf{(5 points)}

\end{enumerate}


\newpage


% ============================================================
% PROBLEM 5
% ============================================================

\begin{problemheader}
\textbf{Problem 5.} \hfill \textbf{15 points}
\end{problemheader}

A factory produces two types of electronic components. Each component is
independently classified as Type I with probability \(0.5\) and Type II
with probability \(0.5\).

Suppose that \(60\) components are produced.

Let \(X\) denote the number of Type I components among the \(60\)
components.

\begin{enumerate}[label=(\alph*),leftmargin=*,itemsep=0.9em]

    \item State the distribution of \(X\).

    Determine whether a normal approximation is appropriate by verifying
    the required conditions. Then state the approximating normal
    distribution, including its mean and variance.

    \hfill \textbf{(6 points)}

    \item Using the normal approximation with an appropriate continuity
    correction, approximate the probability that among the \(60\)
    components there are \textbf{at least \(27\) components of each type}.

    \hfill \textbf{(9 points)}

\end{enumerate}

\end{document}